%%%%%%%%%%%%%%%%
%%% deut 13 %%%
%%%%%%%%%%%%%%%%

\Chapter{13}

\Verse{1}
{
	\hDtn{אֵת כּל הַדָּבָר אֲשֶׁר אָנֹכִי מְצַוֶּה אֶתְכֶם אֹתוֹ תִשְׁמְרוּ לַעֲשׂוֹת לֹא תֹסֵף עָלָיו וְלֹא תִגְרַע מִמֶּנּוּ׃}
}
{
	\eDtn{
		“Everything that I command you: you shall be watchful
		to do it. You shall not add onto it, and you shall not
		subtract from it.
	}
}
{
	Okay so now Deuteronomy, which is Greek for ``second law'' because it's sort
	of like a sequel to the bible, goes on to say that you shouldn't add or
	subtract from the bible.

	This is one commandment that is violated precisely by the authors of the bible
	and no one else.

	This may be at least part of the reason why neither religious people nor
	atheists seem to write bibles very often.

%	\footnote{
%		You shall not add onto it, and you shall not subtract from it. Moses
%		repeats the principle of not adding to or taking away from the law. This
%		repetition emphasizes its extreme importance. It is a command that
%		affects all other commandments. See the comment on Deut 4:2.
%	}
}

\Verse{2}
{
	\hDtn{כִּי יָקוּם בְּקִרְבְּךָ נָבִיא אוֹ חֹלֵם חֲלוֹם וְנָתַן אֵלֶיךָ אוֹת אוֹ מוֹפֵת׃}
}
{
	\eDtn{
		“When a prophet or one who has a dream will get up among
		you and will give you a sign or a wonder,
	}
}
{
	Ok so if any prophets or people who say they had a dream get up and do like a
	miracle or whatever,
}

\Verse{3}
{
	\hDtn{וּבָא הָאוֹת וְהַמּוֹפֵת אֲשֶׁר דִּבֶּר אֵלֶיךָ לֵאמֹר נֵלְכָה אַחֲרֵי אֱלֹהִים אֲחֵרִים אֲשֶׁר לֹא יְדַעְתָּם וְנעבְדֵם׃}
}
{
	\eDtn{
		and the sign or the wonder of which he spoke to
		you—saying, ‘Let’s go after other gods,’ whom you haven’t
		known, ‘and let’s serve them’—will come to pass,
	}
}
{
	and the lesson of the miracle is ``There's a different god aside from Yahweh
	who's actually better at godding.''
}

\Verse{4}
{
	\hDtn{לֹא תִשְׁמַע אֶל דִּבְרֵי הַנָּבִיא הַהוּא אוֹ אֶל חוֹלֵם הַחֲלוֹם הַהוּא כִּי מְנַסֶּה יְהֹוָה אֱלֹהֵיכֶם אֶתְכֶם לָדַעַת הֲיִשְׁכֶם אֹהֲבִים אֶת יְהֹוָה אֱלֹהֵיכֶם בְּכל לְבַבְכֶם וּבְכל נַפְשְׁכֶם׃}
}
{
	\eDtn{
		you shall not listen to that prophet’s words or to that
		one who has the dream, because YHWH, your God, is testing
		you, to know whether you are loving YHWH, your God, with
		all your heart and with all your soul.
	}
}
{
	don't listen, that's just Yahweh testing you to see if you love Yahweh—he's
	your God remember—with all your heart and all your soul.
}

\Verse{5}
{
	\hDtn{אַחֲרֵי יְהֹוָה אֱלֹהֵיכֶם תֵּלֵכוּ וְאֹתוֹ תִירָאוּ וְאֶת מִצְוֺתָיו תִּשְׁמֹרוּ וּבְקֹלוֹ תִשְׁמָעוּ וְאֹתוֹ תַעֲבֹדוּ וּבוֹ תִדְבָּקוּן׃}
}
{
	\eDtn{
		You shall go after YHWH, your God, and you shall fear Him,
		and you shall observe His commandments and listen to His
		voice and serve Him and cling to Him.
	}
}
{
	Go after Yahweh, he's your God.

	And be afraid of him.

	And do what he says and listen to him and cling to him.
}

\Verse{6}
{
	\hDtn{וְהַנָּבִיא הַהוּא אוֹ חֹלֵם הַחֲלוֹם הַהוּא יוּמָת כִּי דִבֶּר סָרָה עַל יְהֹוָה אֱלֹהֵיכֶם הַמּוֹצִיא אֶתְכֶם מֵאֶרֶץ מִצְרַיִם וְהַפֹּדְךָ מִבֵּית עֲבָדִים לְהַדִּיחֲךָ מִן הַדֶּרֶךְ אֲשֶׁר צִוְּךָ יְהֹוָה אֱלֹהֶיךָ לָלֶכֶת בָּהּ וּבִעַרְתָּ הָרָע מִקִּרְבֶּךָ׃}
}
{
	\eDtn{
		And that prophet or that one who has the dream shall be
		put to death, because he spoke a misrepresentation about
		YHWH, your God, who brought you out from the land of Egypt
		and who redeemed you from a house of slaves, to drive you
		from the way in which YHWH, your God, commanded you to go.
		So you shall burn away what is bad from among you.
	}
}
{
	And that person from before who had the dream, we kill 'em!

	Because he's misrepresenting Yahweh.

	What do you mean who's Yahweh?

	He's \emph{your} God, who brought \emph{you} out of the land of Egypt, don't
	worry about the Levites yes I know we have Egyptian names and nobody else does
	just don't—

	...don't worry about that.

	Burn away all that bad stuff from among you.

	That's a good line isn't it?

	I'm gonna use that one NINE TIMES right here in Deuteronomy folks so strap in
	we've got a lot more Deuteronomy to cover.

%	\footnote{
%		burn away what is bad. This is the first of nine occurrences of this
%		expression in Deuteronomy. Until now Moses has used the word literally
%		to describe the burning fire on Mount Sinai (see the comment on 4:11).
%		Now he turns it from describing what God has done to conveying
%		figuratively what the people of Israel must do themselves. This device
%		of using a word in two ways so as to compare what God does and what
%		humans do occurs elsewhere in Deuteronomy as well. Notably, in this very
%		verse, the false prophet or dreamer is described as “driving” Israel
%		from their God. Later God will be described as “driving” Israel from its
%		land because of their apostasy (see the comment on Deut 30:1). In this
%		way, Deuteronomy is an exquisitely woven tapestry of language, pointing
%		to connections and to causes and effects in human behavior.
%	}
}

\Verse{7}
{
	\hDtn{כִּי יְסִיתְךָ אָחִיךָ בֶן אִמֶּךָ אוֹ בִנְךָ אוֹ בִתְּךָ אוֹ אֵשֶׁת חֵיקֶךָ אוֹ רֵעֲךָ אֲשֶׁר כְּנַפְשְׁךָ בַּסֵּתֶר לֵאמֹר נֵלְכָה וְנַעַבְדָה אֱלֹהִים אֲחֵרִים אֲשֶׁר לֹא יָדַעְתָּ אַתָּה וַאֲבֹתֶיךָ׃}
}
{
	\eDtn{
		“When your brother, your mother’s son or your father’s
		son, or the wife of your bosom, or your friend who is as
		your own self will entice you in secret, saying, ‘Let’s go
		and serve other gods,’ whom you haven’t known, you and
		your fathers,
	}
}
{
	Ok so, when your family (I'll briefly list all the different types of family
	member) comes and whispers in your ear saying ``Hey I've been seeing this hot
	new god and he doesn't live in a tent like your god he really has Ba'als''—some
	random new guy who you never met and your dad hasn't either,
}

\Verse{8}
{
	\hDtn{מֵאֱלֹהֵי הָעַמִּים אֲשֶׁר סְבִיבֹתֵיכֶם הַקְּרֹבִים אֵלֶיךָ אוֹ הָרְחֹקִים מִמֶּךָּ מִקְצֵה הָאָרֶץ וְעַד קְצֵה הָאָרֶץ׃}
}
{
	\eDtn{
		from the gods of the peoples who are all around you, those
		close to you or those far from you, from one end of the
		earth to the other end of the earth,
	}
}
{
	from the gods of the natives, wherever they are,
}

\Verse{9}
{
	\hDtn{לֹא תֹאבֶה לוֹ וְלֹא תִשְׁמַע אֵלָיו וְלֹא תָחוֹס עֵינְךָ עָלָיו וְלֹא תַחְמֹל וְלֹא תְכַסֶּה עָלָיו׃}
}
{
	\eDtn{
		you shall not consent to him, and you shall not listen to
		him, and your eye shall not pity him, and you shall not
		have compassion and shall not cover it up for him.
	}
}
{
	don't do it, don't listen, and don't feel pity or compassion and don't keep
	those secrets he was whispering.
}

\Verse{10}
{
	\hDtn{כִּי הָרֹג תַּהַרְגֶנּוּ יָדְךָ תִּהְיֶה בּוֹ בָרִאשׁוֹנָה לַהֲמִיתוֹ וְיַד כּל הָעָם בָּאַחֲרֹנָה׃}
}
{
	\eDtn{
		But you shall kill him. Your hand shall be on him first to
		put him to death, and all the people’s hand thereafter.
	}
}
{
	KILL HIM.

	Yes you specifically.

	Kill your friend or whoever it was.

	You start killing him and then everyone else in town joins in.
}

\Verse{11}
{
	\hDtn{וּסְקַלְתּוֹ בָאֲבָנִים וָמֵת כִּי בִקֵּשׁ לְהַדִּיחֲךָ מֵעַל יְהֹוָה אֱלֹהֶיךָ הַמּוֹצִיאֲךָ מֵאֶרֶץ מִצְרַיִם מִבֵּית עֲבָדִים׃}
}
{
	\eDtn{
		And you shall stone him with stones so he dies because he
		sought to drive you away from YHWH, your God, who brought
		you out from the land of Egypt, from a house of slaves.
	}
}
{
	Hit him with rocks until he dies because he tried to take you away from
	Yahweh.

	But that's crazy! You know Yahweh, he's your God, who brought you out of the
	land of Egypt, remember Egypt?

	No?

	Ok moving on.
}

\Verse{12}
{
	\hDtn{וְכל יִשְׂרָאֵל יִשְׁמְעוּ וְיִרָאוּן וְלֹא יוֹסִפוּ לַעֲשׂוֹת כַּדָּבָר הָרָע הַזֶּה בְּקִרְבֶּךָ׃}
}
{
	\eDtn{
		And all Israel will hear and fear and won’t continue to do
		a bad thing like this among you.
	}
}
{
	Then everyone will be too scared to ever do bad stuff like that again.
}

\Verse{13}
{
	\hDtn{כִּי תִשְׁמַע בְּאַחַת עָרֶיךָ אֲשֶׁר יְהֹוָה אֱלֹהֶיךָ נֹתֵן לְךָ לָשֶׁבֶת שָׁם לֵאמֹר׃}
}
{
	\eDtn{
		“When you’ll hear in one of your cities that YHWH, your
		God, is giving you to live there, saying,
	}
}
{
	Ok now, if you're in a city, one of the ones Yahweh, \emph{your} God, is giving
	you, and you hear someone saying,
}

\Verse{14}
{
	\hDtn{יָצְאוּ אֲנָשִׁים בְּנֵי בְלִיַּעַל מִקִּרְבֶּךָ וַיַּדִּיחוּ אֶת יֹשְׁבֵי עִירָם לֵאמֹר נֵלְכָה וְנַעַבְדָה אֱלֹהִים אֲחֵרִים אֲשֶׁר לֹא יְדַעְתֶּם׃}
}
{
	\eDtn{
		‘Good-for-nothing people have gone out from among you and
		driven away their city’s residents, saying: “Let’s go and
		serve other gods,” whom you haven’t known,’
	}
}
{
	the same thing as before basically,

%	\footnote{
%		Good-for-nothing. The Hebrew blîya‘al appears to be a compound of
%		two words meaning “without value” or “without benefit.” This corresponds
%		to the English expression “good-for-nothing,” which fits with every
%		occurrence of the expression in the Tanak’s narrative books.
%	}
}

\Verse{15}
{
	\hDtn{וְדָרַשְׁתָּ וְחָקַרְתָּ וְשָׁאַלְתָּ הֵיטֵב וְהִנֵּה אֱמֶת נָכוֹן הַדָּבָר נֶעֶשְׂתָה הַתּוֹעֵבָה הַזֹּאת בְּקִרְבֶּךָ׃}
}
{
	\eDtn{
		and you’ll inquire and investigate and ask well, and,
		here, the thing is true—this offensive thing was done
		among you—
	}
}
{
	and you'll look into it and realize it's true—this offensive thing happened
	right here among you—
}

\Verse{16}
{
	\hDtn{הַכֵּה תַכֶּה אֶת יֹשְׁבֵי הָעִיר הַהִוא לְפִי חָרֶב הַחֲרֵם אֹתָהּ וְאֶת כּל אֲשֶׁר בָּהּ וְאֶת בְּהֶמְתָּהּ לְפִי חָרֶב׃}
}
{
	\eDtn{
		you shall strike that city’s residents by the sword,
		completely destroying it and everything in it and its
		animals with the sword.
	}
}
{
	cut 'em all up into little pieces and kill all the people in the city and all
	the animals too with the sword.

%	\footnote{
%		completely destroying it and everything in it. On the extreme violence
%		that is commanded in destroying those (Israelites) who would lead one to
%		worship other gods, or in destroying an entire Israelite town that
%		worships other gods (including any righteous; cf. Sodom): again the
%		Torah distinguishes between ethical sins and ritual sins. Ethical sins
%		(those that occur between a human and one’s fellow humans; theft, for
%		example) are judged in terms of each individual person’s actions, and
%		the person’s intent is taken into account. Ritual sins (those that are
%		between a human and God; blasphemy, for example) are judged in terms of
%		groups and zones, and intent is not an issue. Thus, when Achan violates
%		the law of erem in Joshua 7, the text says that “the children of
%		Israel” violated the erem (7:1)—and all of Israel suffers a
%		defeat. Achan takes items that were forbidden to the Israelites, and he
%		hides them in his tent; and the result is that not only is Achan stoned
%		and burned, but his family, his animals, everything that was in the
%		tent, and also the tent itself are stoned and burned. Questions about
%		the harsh response to Israelites who turn to (or seduce others to) pagan
%		worship must be addressed in this light. Worship of other gods is the
%		ultimate ritual sin in Deuteronomy. As with all ritual infringements,
%		the perpetrators must be utterly eliminated from Israel. If it is an
%		entire city, then that city must be eliminated from Israel—including any
%		righteous members of the community and even the animals. All of this may
%		seem foreign and barbaric in our time. That is because (1) Israel has
%		not been a religious state for 2000 years; and (2) after the destruction
%		of the Second Temple, ritual became less central in Judaism, and
%		emphasis on ethics rose. Also, (3) monotheism has defeated pagan
%		religion in Western civilization, so we no longer feel the degree of
%		threat that authors of the Bible felt toward it.
%	}
}

\Verse{17}
{
	\hDtn{וְאֶת כּל שְׁלָלָהּ תִּקְבֹּץ אֶל תּוֹךְ רְחֹבָהּ וְשָׂרַפְתָּ בָאֵשׁ אֶת הָעִיר וְאֶת כּל שְׁלָלָהּ כָּלִיל לַיהֹוָה אֱלֹהֶיךָ וְהָיְתָה תֵּל עוֹלָם לֹא תִבָּנֶה עוֹד׃}
}
{
	\eDtn{
		And you shall gather all its spoil into the middle of its
		square, and you shall burn the city and all its spoil in
		fire entirely to YHWH, your God, and it shall be a tell
		eternally. It shall not be rebuilt again.
	}
}
{
	Then you gather all their stuff up into a big pile, and burn the city and all
	its stuff in fire to Yahweh, \emph{your} God, and it'll be a big pile forever
	so everyone knows never to rebuild the city again!

%	\footnote{
%		a tell. A tell is a hill that came to exist as a result of successive
%		layers of cities being built on top of one another. Because a tell is
%		composed of sand and rock and debris rather than soil, very little grass
%		or vegetation grows on it. Such a mound is a treasure to an
%		archaeologist, but it is a curse here in the text: not only is an
%		apostate city to be destroyed, but it is to become a tell permanently,
%		never to be rebuilt.
%	}
}

\Verse{18}
{
	\hDtn{וְלֹא יִדְבַּק בְּיָדְךָ מְאוּמָה מִן הַחֵרֶם לְמַעַן יָשׁוּב יְהֹוָה מֵחֲרוֹן אַפּוֹ וְנָתַן לְךָ רַחֲמִים וְרִחַמְךָ וְהִרְבֶּךָ כַּאֲשֶׁר נִשְׁבַּע לַאֲבֹתֶיךָ׃}
}
{
	\eDtn{
		And let nothing from the complete destruction cling in
		your hand, so that YHWH will turn back from His flaring
		anger and will give you mercy. And He’ll be merciful to
		you and multiply you as He swore to your fathers
	}
}
{
	Don't keep anything. Burn it all.

	That way Yahweh won't hit you with that burning anger of his and he'll give
	you mercy instead.

	And he'll multiply your seed like the stars in the sky and give you more
	p\eR{aternit}y than you can possibly f\eR{ather} in a lifetime!
}

\Verse{19}
{
	\hDtn{כִּי תִשְׁמַע בְּקוֹל יְהֹוָה אֱלֹהֶיךָ לִשְׁמֹר אֶת כּל מִצְוֺתָיו אֲשֶׁר אָנֹכִי מְצַוְּךָ הַיּוֹם לַעֲשׂוֹת הַיָּשָׁר בְּעֵינֵי יְהֹוָה אֱלֹהֶיךָ׃}
}
{
	\eDtn{
		when you’ll listen to the voice of YHWH, your God, to
		observe all His commandments that I command you today, to
		do what is right in the eyes of YHWH, your God.
	}
}
{
	Do it.

	You know it's right.

	End of chapter!
}
